\section{Introduction}
\label{sec:introduction}

Driving simulation and scene editing require 3D scenes that support changes in camera viewpoint, vehicle motion, and human pose~\cite{yang2023unisim,zhou2024hugsim}. However, driving images reveal only the surfaces visible from the recorded viewpoints and actor configurations~\cite{sun2020waymo,caesar2020nuscenes}. Moving a camera, turning a vehicle, or changing a pedestrian's pose exposes previously unobserved regions. Reconstructing scenes for these operations therefore requires completing the missing geometry and appearance while ensuring that completed vehicle and human surfaces move with their respective actors.

Recent feed-forward methods fuse multiple images to directly predict driving-scene geometry and Gaussian representations~\cite{yang2025storm,chen2025dggt,yu2026streetforward,chi2026pointforward}. Motion estimation and visibility modeling further help align dynamic objects and integrate observations over time. Multi-frame aggregation increases coverage of surfaces visible in the input sequence. When a new viewpoint or actor motion reveals surfaces unseen throughout the sequence, reconstruction must infer their geometry and appearance.

Generative reconstruction uses asset priors, Gaussian completion, or generated views to infer unobserved content~\cite{yang2025genassets,liao2025complete,streetcrafter,zhu2025worldsplat}. Meanwhile, compositional reconstruction represents the background and moving actors separately to support vehicle motion and human pose control~\cite{yan2024street,chen2025omnire,zhou2024hugsim}. Applying generative completion to such scenes requires generated surfaces to move with their actors. When an actor is only partially observed, information from the background and other actors can provide additional context for inferring its missing content.

To this end, we introduce DriveCompose, a framework that generatively reconstructs compositional driving scenes from image sequences while supporting independent actor control. We represent the background in world coordinates, vehicle surfaces in each vehicle's own coordinate frame, and human surfaces in a reference pose with skeletal bindings~\cite{yan2024street,chen2025omnire}. Completion operates directly on these representations, allowing generated vehicle and human surfaces to follow their actors' motion and pose changes (Sec.~\ref{sec:representation}). The model jointly completes missing background and actor content using shared scene context.

As illustrated in Fig.~\ref{fig:overview}, DriveCompose first recovers scene geometry and actor motion from the image sequence and aligns multi-frame image features with the background, vehicle, and human representations. The completion model captures local structure within each component, while cross-component memory gathers information from the background and actors and distributes the shared context back to each component. Missing regions of a vehicle or human can therefore be completed using information from the surrounding scene. The model learns to generate missing latent codes through conditional flow matching~\cite{lipman2022flowmatching}, and component-specific decoders convert the completed representations into a Gaussian scene~\cite{kerbl2023gaussians}.

To maintain consistency with input observations, we preserve the geometry and appearance of reliably observed surfaces during completion. We then render the completed scene at the recorded camera viewpoints and actor states and compare it with the input images to check for inconsistencies introduced by completion. Scenes that pass this check can reuse the same completed content for rendering under different camera viewpoints and actor states (Sec.~\ref{sec:composition}).

On Waymo, DriveCompose improves dynamic scene reconstruction and off-trajectory view synthesis over the evaluated baselines (Sec.~\ref{sec:experiments}). Component ablations further show that cross-component information sharing reduces mean absolute RGB error on unseen surfaces, with larger absolute error reductions under sparser observations (Sec.~\ref{sec:ablations}).

\Needspace{6\baselineskip}
Our contributions are:
\begin{itemize}
\item We introduce DriveCompose, a generative reconstruction framework for closed-loop driving simulation. It recovers and completes backgrounds, vehicles, and humans from image sequences to construct compositional Gaussian scene assets. The resulting scene reuses the same completed assets for composition and rendering as camera poses, vehicle poses, and human poses change.
\item We develop a cross-component latent completion model based on conditional flow matching. Local attention models spatial relationships within each component, while cross-component memory aggregates background and actor features to provide scene context for generating each component's missing content.
\item We evaluate the framework through dynamic reconstruction and off-trajectory view synthesis, and use component ablations to quantify the reduction in mean absolute RGB error on unseen surfaces from cross-component information sharing.
\end{itemize}

Section~\ref{sec:related} reviews related work, Section~\ref{sec:method} presents the method, and Section~\ref{sec:experiments} reports experiments. Sections~\ref{sec:discussion} and~\ref{sec:conclusion} discuss limitations and conclude.
